% =====================================================================
% UoN Beamer presentation — personal working template
% Based on the University of Nottingham Beamer theme (UoN.sty),
% itself derived from the public UoN/UNNC Beamer template.
% Engine : XeLaTeX
% Build  : latexmk -xelatex main.tex
% =====================================================================
\documentclass[10pt,aspectratio=169]{beamer}
\usepackage{UoN}

% ---- Title / cover page: edit titlepage.tex (kept separate for reusers) ----
% =====================================================================
% Title (cover) page — edit the fields below, then % =====================================================================
% Title (cover) page — edit the fields below, then % =====================================================================
% Title (cover) page — edit the fields below, then \input{titlepage}
% from main.tex. Kept separate so anyone reusing the template only has
% to change this one file. Delete the fields you don't need.
% =====================================================================
\title{Nottingham Presentation Template}
\subtitle{A reusable Beamer deck for talks and seminars}
\author{Your Name}
\institute{University of Nottingham}
\date{\today}

\begin{document}

\maketitle

% from main.tex. Kept separate so anyone reusing the template only has
% to change this one file. Delete the fields you don't need.
% =====================================================================
\title{Nottingham Presentation Template}
\subtitle{A reusable Beamer deck for talks and seminars}
\author{Your Name}
\institute{University of Nottingham}
\date{\today}

\begin{document}

\maketitle

% from main.tex. Kept separate so anyone reusing the template only has
% to change this one file. Delete the fields you don't need.
% =====================================================================
\title{Nottingham Presentation Template}
\subtitle{A reusable Beamer deck for talks and seminars}
\author{Your Name}
\institute{University of Nottingham}
\date{\today}

\begin{document}

\maketitle


% ---------------------------------------------------------------
\begin{frame}{Contents}
  \tableofcontents
\end{frame}

% =================================================================
\section{Introduction}

\begin{frame}{Why this template}
  \begin{itemize}[<+->]
    \item The \texttt{UoN.sty} theme already encodes the university's
          visual identity --- colours, logo header, progress bar.
    \item You write content; \LaTeX\ handles consistent typesetting.
    \item Produces a portable \texttt{pdf} --- no font or version drift.
    \item Fonts come from TeX Live itself (\texttt{noto} + Fandol via
          \texttt{ctex}), so the build never depends on system fonts.
  \end{itemize}
\end{frame}

\begin{frame}{The brand palette}
  The theme uses the official University of Nottingham colours:
  \begin{columns}[T]
    \begin{column}{0.48\textwidth}
      \begin{itemize}
        \item Nottingham Blue
              \textcolor{NottinghamBlue}{\rule{1.2em}{1em}}
              \texttt{\#10263B} --- primary
        \item Sky Blue
              \textcolor{MalaysiaSkyBlue}{\rule{1.2em}{1em}}
              \texttt{\#009BC1} --- accent
      \end{itemize}
    \end{column}
    \begin{column}{0.48\textwidth}
      \begin{block}{Body text}
        Normal text is set in Nottingham Blue on white --- never pure
        black --- for a softer, official look.
      \end{block}
    \end{column}
  \end{columns}
\end{frame}

% =================================================================
\section{Content building blocks}

\begin{frame}{Columns and structure}
  \begin{columns}[T]
    \begin{column}{0.58\textwidth}
      \begin{itemize}
        \item Split any frame into columns.
        \item Good for text beside a figure.
        \item Align with \texttt{[T]} or \texttt{[c]}.
      \end{itemize}
    \end{column}
    \begin{column}{0.38\textwidth}
      \begin{alertblock}{Note}
        Keep each slide to one idea.
      \end{alertblock}
    \end{column}
  \end{columns}
\end{frame}

\begin{frame}{Math typesetting}
  \TeX\ is the gold standard for equations:
  \begin{equation*}
    \mathrm{i}\,\hslash\frac{\partial}{\partial t}\Psi(\mathbf{r},t)
    = -\frac{\hslash^2}{2\,m}\nabla^{2}\Psi(\mathbf{r},t)
      + V(\mathbf{r})\Psi(\mathbf{r},t)
  \end{equation*}
  \begin{itemize}
    \item Inline maths such as $E = mc^{2}$ flows with the text.
    \item Numbered environments work too:
          \begin{equation}
            \sum_{i=1}^{n} i = \frac{n(n+1)}{2}
          \end{equation}
  \end{itemize}
\end{frame}

\begin{frame}[fragile]{Code listing}
  \begin{block}{A minimal document}
    \begin{lstlisting}[language=TeX]
\documentclass[10pt,aspectratio=169]{beamer}
\usepackage{UoN}
\begin{document}
\maketitle
\end{document}
    \end{lstlisting}
  \end{block}
\end{frame}

\begin{frame}{Table}
  \begin{table}
    \centering
    \begin{tabular}{lrrr}
      \toprule
      \textbf{Quarter} & \textbf{Revenue} & \textbf{Cost} & \textbf{Growth} \\
      \midrule
      Q1 & 12 & 8 & 4\% \\
      Q2 & 19 & 11 & 7\% \\
      Q3 & 15 & 9 & 5\% \\
      Q4 & 22 & 12 & 9\% \\
      \bottomrule
    \end{tabular}
    \caption{Sample financial summary.}
  \end{table}
\end{frame}

\begin{frame}{A simple figure}
  \centering
  \begin{tikzpicture}
    \draw[NottinghamBlue, very thick, ->] (0,0) -- (6,0) node[right] {$x$};
    \draw[NottinghamBlue, very thick, ->] (0,0) -- (0,3.5) node[above] {$y$};
    \draw[MalaysiaSkyBlue, ultra thick, domain=0:5.5, samples=60]
      plot (\x, {2.2*exp(-0.25*\x)});
    \draw[NottinghamBlue, dashed] (0,2.2) -- (5.5,0.7);
  \end{tikzpicture}

  \vspace{0.5em}
  \footnotesize Figures are drawn with \texttt{tikz}, bundled with the theme.
\end{frame}

% =================================================================
\section{Closing}

\begin{frame}{Key takeaways}
  \begin{enumerate}
    \item Edit the metadata block at the top of \texttt{main.tex}.
    \item Reuse the blocks here when building a real talk.
    \item Build with \texttt{latexmk -xelatex main.tex}.
  \end{enumerate}
\end{frame}

\lastframe

\end{document}
